% Generated; original Preface physical pp. 13--14.
\chapter*{Preface}
\addcontentsline{toc}{chapter}{Preface}
\hypertarget{orig-front-xiii}{}
\textit{Complex Analysis} has successfully maintained its place as the standard elementary text on functions of one complex variable. There is, nevertheless, need for a new edition, partly because of changes in current mathematical terminology, partly because of differences in student preparedness and aims.

There are no radical innovations in the new edition. The author still believes strongly in a geometric approach to the basics, and for this reason the introductory chapters are virtually unchanged. In a few places, throughout the book, it was desirable to clarify certain points that experience has shown to have been a source of possible misunderstanding or difficulties. Misprints and minor errors that have come to my attention have been corrected. Otherwise, the main differences between the second and third editions can be summarized as follows:

\par\hangindent=1.25em\hangafter=1 1. Notations and terminology have been modernized, but it did not seem necessary to change the style in any significant way.

\par\hangindent=1.25em\hangafter=1 2. In Chapter 2 a brief section on the change of length and area under conformal mapping has been added. To some degree this infringes on the otherwise self-contained exposition, for it forces the reader to fall back on calculus for the definition and manipulation of double integrals. The disadvantage is minor.

\par\hangindent=1.25em\hangafter=1 3. In Chapter 4 there is a new and simpler proof of the general form of Cauchy’s theorem. It is due to A. F. Beardon\index{Beardon A F@Beardon, A. F.}, who has kindly permitted me to reproduce it. It complements but does not replace the old proof, which has been retained and improved.

\par\hangindent=1.25em\hangafter=1 4. A short section on the Riemann zeta function has been included.

This always fascinates students, and the proof of the functional equation illustrates the use of residues in a less trivial situation than the mere computation of definite integrals.

\par\hangindent=1.25em\hangafter=1 5. Large parts of Chapter 8 have been completely rewritten. The main purpose was to introduce the reader to the terminology of germs and sheaves while emphasizing all the classical concepts. It goes without saying that nothing beyond the basic notions of sheaf theory would have been compatible with the elementary nature of the book.

\par\hangindent=1.25em\hangafter=1 6. The author has successfully resisted the temptation to include Riemann surfaces as one-dimensional complex manifolds. The book would lose much of its usefulness if it went beyond its purpose of being no more than an introduction to the basic methods and results of complex function theory in the plane.

It is my pleasant duty to thank the many who have helped me by pointing out misprints, weaknesses, and errors in the second edition. I am particularly grateful to my colleague Lynn Loomis, who kindly let me share student reaction to a recent course based on my book.

\begin{flushright}\textit{Lars V. Ahlfors}\end{flushright}

% Original index cross-reference (source physical page 341).
\index{Analytic function|see{Function, analytic}}
